% =========================================================
% 2.5 SAFEROUTE
% El \section{} de esta sección está en el chapter.tex del capítulo;
% sus referencias van en seccion2_5.bib (misma carpeta).
% =========================================================

Cerrando este recorrido por los antecedentes, SafeRoute introduce un enfoque basado en aprendizaje por refuerzo, en contraste con los métodos deterministas (Dijkstra) y estadísticos (regresión logística) revisados anteriormente. Propuesto en 2020, SafeRoute utiliza aprendizaje por refuerzo profundo para la generación de rutas seguras y óptimas en las ciudades de Boston, Nueva York y San Francisco. Sus resultados mostraron una mejora del 17\% en la distancia promedio entre la ruta generada y los sitios de criminalidad, a cambio de un incremento del 7\% en la longitud de la ruta\cite{levy2020saferoute}.

Comparte mucha similitud con el trabajo aplicado en la Ciudad de México, en cuanto ambos se basan en reportes oficiales de crímenes de sus respectivas áreas de estudio como motor para la generación de los modelos. Asimismo, para el presente proyecto, lo más relevante es la definición precisa de los algoritmos de entrenamiento y prueba del agente para la generación de rutas seguras, los cuales podrán retomarse posteriormente en la fase de construcción del modelo final de ruteo de nuestro aplicativo.
